\documentclass[12pt,a4paper]{extarticle} % Используем расширенный класс

\usepackage[T2A]{fontenc}
\usepackage[utf8]{inputenc}
\usepackage[english, russian]{babel}

% Улучшенный Times-образный шрифт
\usepackage{tempora}
\usepackage{newtxmath}
\usepackage{microtype} % Чтобы текст ложился ровно и красиво
\usepackage{type1ec}

\usepackage[
left=30mm,
right=10mm,
top=20mm,
bottom=20mm,
footskip=10mm
]{geometry}

\usepackage[nodisplayskipstretch]{setspace}
\onehalfspacing 

\usepackage{indentfirst}
\setlength{\parindent}{1.25cm}

% Настройка заголовков: центрирование, жирный
\usepackage{titlesec}
\titleformat{\section}{\normalfont\Large\bfseries\centering}{}{0pt}{}
\titleformat{\subsection}{\normalfont\large\bfseries\centering}{}{0pt}{}

\usepackage{graphicx}
\usepackage{caption}

% Нумерация страниц в центре нижнего поля
\usepackage{fancyhdr}
\pagestyle{fancy}
\fancyhf{}
\fancyfoot[C]{\thepage}
\renewcommand{\headrulewidth}{0pt}

%Правильные сниппеты с кодом
\usepackage{minted}
\setminted{
breaklines=true,
breakanywhere=true,
fontsize=\footnotesize,
baselinestretch=1.0,  % Устанавливает одинарный интервал внутри кода
breaksymbol={},           % Удаляем символ переноса
fontfamily=tt,
xleftmargin=-2.15cm,  % Делает отступ как у красной строки
}
% Отступы у списков
\usepackage{enumitem}

% Настройка для всех маркированных списков
\setlist[itemize]{
noitemsep,           % Убирает интервал между пунктами
topsep=0pt,          % Убирает отступ перед списком
parsep=0pt,          % Убирает отступ внутри пункта (если текст в несколько строк)
partopsep=0pt,
leftmargin=1.25cm    % Выравнивание по красной строке
}

% Настройка для нумерованных списков
\setlist[enumerate]{
noitemsep,
topsep=4pt,
leftmargin=1.25cm
}
% Настройка гиперссылок в документе
\usepackage[unicode, pdftex]{hyperref}
\usepackage{url}

% Форматирование приложений
\usepackage[russian]{cleveref}
\crefformat{appsec}{приложении~#2#1#3}
\Crefformat{appsec}{Приложении~#2#1#3}
\makeatother

\hypersetup{
citecolor=green,       % Цвет ссылок на литературу
filecolor=blue,     
urlcolor=blue,         % Цвет внешних URL-ссылок
pdftitle={Разработка веб-сайта для нахождения оптимальной рассадки учеников в классе}, % Метаданные PDF
pdfauthor={Прокофьев Даниил}
pdfcopyright={CC BY-SA 4.0}
}

% Картинки
\usepackage{graphicx}
\usepackage{float}
\graphicspath{{../presentation/public/screenshots}}

% Иконка лицензии CC-BY-SA 4.0
\usepackage{ccicons}

% =================================================================
\begin{document}

	% --- ТИТУЛЬНЫЙ ЛИСТ ---
	\noindent
	\begin{minipage}[b]{0.5\textwidth}
		{\begingroup
			
			\color{gray}
			\footnotesize % Уменьшаем размер
			\ccbysa \\ % Векторный значок CC BY-SA
			Данная работа доступна по лицензии \\
			\href{https://creativecommons.org/licenses/by-sa/4.0/deed.ru}{Creative Commons Attribution-ShareAlike 4.0 International}
			
			\endgroup}
	\end{minipage}%
	\begin{minipage}[b]{0.5\textwidth}
		\raggedleft
		\textbf{Прокофьев Даниил Вячеславович} \\
		МБОУ «Гимназия» г. Обнинска, 10 класс \\
		Научный руководитель: Утянская Е. В.
	\end{minipage}
	
	\vspace{2cm}
	
	\begin{center}
		Секция: Информатика \\[1cm]
		{\Large \textbf{Разработка веб-сайта для нахождения оптимальной рассадки учеников в классе}}
	\end{center}
	
	Эффективная организация образовательной среды напрямую влияет на сохранение здоровья обучающихся и качество усвоения материала. Традиционный метод рассадки учащихся «вручную» зачастую субъективен и не позволяет в полной мере учесть совокупность факторов: рост ученика, состояние зрения и слуха, а также социальное взаимодействие. В связи с этим разработка автоматизированной системы, позволяющей создать рассадку, учитывающую все ограничения и предпочтения, является актуальной задачей.
	
	\textbf{Цель работы:} создать веб-сайт, позволяющий учителю в автоматизированном режиме рассадить учеников в классе согласно медицинским ограничениям и социальным связям между учениками.
	
	В работе решаются задачи:
	
	\begin{enumerate}[noitemsep]
		\item Выбор подходящих технологий для создания веб-сайта
		\item Выбор, написание, отладка алгоритма для генерации оптимальной рассадки с заданными ограничениями
		\item Автоматизация сборки компонентов веб-сайта и предоставление каждому возможности развернуть свой собственный экземпляр генератора
		\item Обеспечение доступа к сайту
	\end{enumerate}
	
	\textbf{Краткие выводы}
	
	\begin{enumerate}
		\item В качестве технологий для создания сайта были выбраны \textbf{язык программирования Go} — для бэкенда и \textbf{JavaScript} с фреймворком \textbf{Vue.js} — для фронтенда
		\item В ходе решения данной задачи был написан сервер на языке программирования \textbf{Go}, использующий для решения поставленной задачи \textbf{меметический алгоритм}. Входные данные и результат работы передаются в формате JSON. Вместе с итоговой рассадкой возвращается не только расположение учеников, но и удовлетворенность рассадкой каждого из них, что позволяет впоследствии отображать тепловую карту учтенности предпочтений на фронтенде
		\item В ходе выполнения данной задачи создан \textbf{мета-репозиторий} проекта на GitHub, содержащий инструкции и Docker Compose файл, позволяющий запустить генератор на любом компьютере с установленным Docker и Docker Compose. Проект разделен на два репозитория — один из них содержит фронтенд, а второй — бэкенд. Это позволило автоматизировать сборку образов Docker контейнеров при помощи GitHub Actions
		\item Для обеспечения доступа к сайту был использован VPS сервер, арендованный у одного из крупных российских провайдеров. Быстрому и удобному развертыванию способствуют полученные в ходе решения 3-ей задачи наработки
	\end{enumerate}
	
	Тестирование программы показало, что если оптимальная рассадка достижима, то она и будет результатом работы программы. В противном случае алгоритм будет отдавать приоритет тем ограничениям, что имеют наибольший приоритет согласно настройкам пользователя.
	
	\begin{figure}[H]
		\centering
		\includegraphics[width=0.8\textwidth]{generator.png}
		\caption{Интерфейс страницы генерации} 
		\label{fig:generator}
	\end{figure}
	
	Весь исходный код доступен по лицензии GNU AGPLv3 и открыт для улучшения любому желающему. Подробную информацию можно получить в мета-репозитории проекта \cite{github_repo}
	
	\bibliographystyle{plainurl}
	\bibliography{references}
\end{document}